\section{Conclusion}
\label{sec:conclusion}

In this work, we introduced \textbf{2D-STEM}, a training-free 2D Spatio-Temporal Exponential Moving Average for stateful dynamic model merging at the edge. By applying Occam's razor to stateful serving, we demonstrated that highly parameterized dynamical frameworks (such as non-equilibrium biochemical kinetics or learned PAC-Bayesian recurrences) can be significantly simplified. Instead, a simple, mathematically pure 2D bilinear recursive filter achieves state-of-the-art results. 

Through formal analysis, we proved that 2D-STEM is analytically guaranteed to preserve the probability simplex at all layers and steps without the need for computational projection or re-normalization operations. To resolve the fundamental transition lag inherent in stateful temporal systems, we introduced Adaptive Temporal Gating (ATG), which detects task switches on-the-fly and instantly resets sequence-level memory to enable seamless task-switching. Our extensive evaluations on the Analytical Coordinate Sandbox show that 2D-STEM reduces excess routing noise above the physical boundary limit by up to $15.8\times$ (and absolute routing jitter by up to $2.75\times$), while outperforming the constant-inertia ChemMerge proxy by up to a massive $51.88\%$ absolute accuracy under rapidly switching heterogeneous streams—all with zero parameter or optimization overhead.

Our findings serve as a broader call to the machine learning community: as we scale models and serving systems, we should prioritize the study of simple, robust, and interpretable baselines before adopting highly parameterized frameworks. Relying on clean, mathematically rigorous formulations often yields serving systems that are not only more transparent, but highly robust and performant. 

In future work, we plan to provide a concrete, physical validation of 2D-STEM on large-scale architectures. Specifically, we will fine-tune a set of specialized LoRA experts (rank $r = 8$ or $16$) on a Vision Transformer (ViT-Base) backbone across diverse downstream tasks (such as CIFAR-100 and DomainNet sub-domains). We will deploy this multi-expert model in dynamic streaming environments to verify that the coordinate-space representational dynamics analyzed in our sandbox environment generalize seamlessly to physical weights. Beyond visual classification, we also plan to extend 2D spatio-temporal bilinear filtering to token-level routing in sparse Mixture-of-Experts (MoE). Specifically, token-level routing suffers from severe representational jitter across adjacent layers and consecutive tokens in a sequence. Applying 2D-STEM to MoE networks would involve propagating routing coefficients across both network layers (spatial) and token sequences (temporal). To handle syntactic boundaries (such as punctuation tokens or special end-of-sentence delimiters), we can implement a syntactic boundary gating mechanism that dynamically scales down the temporal momentum, resetting the state propagation when transitioning between semantic segments. This would stabilize token-level expert assignments, preventing routing thrashing, and reducing latency on autoregressive edge-generation backbones. Finally, we plan to investigate dynamic cross-attention parameter blending in multi-modal generative models.
